\documentclass[11pt,a4paper]{article}

\usepackage[utf8]{inputenc}
\usepackage[T1]{fontenc}
\usepackage{amsmath,amssymb,amsthm,mathrsfs}
\usepackage{geometry}
\usepackage{hyperref}
\usepackage{enumitem}
\usepackage{booktabs}
\usepackage{graphicx}
\usepackage{xcolor}
\usepackage{microtype}

\geometry{margin=2.5cm}

\hypersetup{
  colorlinks=true,
  linkcolor=blue!60!black,
  citecolor=blue!60!black,
  urlcolor=blue!60!black
}

\newtheorem{theorem}{Theorem}[section]
\newtheorem{proposition}[theorem]{Proposition}
\newtheorem{lemma}[theorem]{Lemma}
\newtheorem{corollary}[theorem]{Corollary}
\newtheorem{conjecture}[theorem]{Conjecture}
\theoremstyle{definition}
\newtheorem{definition}[theorem]{Definition}
\newtheorem{remark}[theorem]{Remark}
\newtheorem{example}[theorem]{Example}

\DeclareMathOperator{\PSL}{PSL}
\DeclareMathOperator{\SL}{SL}
\DeclareMathOperator{\GL}{GL}
\DeclareMathOperator{\SU}{SU}
\DeclareMathOperator{\Cl}{Cl}
\DeclareMathOperator{\Gal}{Gal}
\DeclareMathOperator{\Reg}{Reg}
\DeclareMathOperator{\Vol}{Vol}
\DeclareMathOperator{\Tr}{Tr}
\DeclareMathOperator{\disc}{disc}
\DeclareMathOperator{\rk}{rk}
\DeclareMathOperator{\Sel}{Sel}

\newcommand{\OK}{\mathcal{O}_K}
\newcommand{\HH}{\mathbb{H}}
\newcommand{\RR}{\mathbb{R}}
\newcommand{\CC}{\mathbb{C}}
\newcommand{\QQ}{\mathbb{Q}}
\newcommand{\ZZ}{\mathbb{Z}}
\newcommand{\NN}{\mathbb{N}}
\newcommand{\FF}{\mathbb{F}}

\title{The ECI Framework Applied to Clay Millennium Problems:\\
A Survey of Strong, Mixed, and Negative Bridges\\
(with honest scope qualifications)}
\author{K\'evin R\'emondi\`ere\\
\small Independent Researcher, Oloron-Sainte-Marie, France\\
\small ORCID: \href{https://orcid.org/0009-0008-2443-7166}{0009-0008-2443-7166}\\
\small \texttt{kevin.remondiere@gmail.com}}
\date{19 May 2026}

\begin{document}

\maketitle

\begin{abstract}
We survey the \emph{Empirical Cross-Galois Identities} (ECI) framework
and its current reach across the seven Clay Millennium Problems. ECI
organizes arithmetic data on imaginary quadratic fields~$K$ through
hyperbolic Bianchi 3-orbifolds $X_K = \HH^{3}/\PSL_{2}(\OK)$ and a small
``Bianchi--ASP'' dictionary that pairs spectral, automorphic and algebro--geometric
objects on $X_{K}$ with Clay-style invariants on number-theoretic targets.
Our goals are deliberately modest. ECI is \textbf{not} a master proof and
\textbf{not} a theory of everything; it is a working organizing principle
supported by~$\sim 70$ verified arithmetic anchors and falsified at
several explicit places. We classify the seven problems into three
honest tiers:
\textbf{(MIXED / CONDITIONAL)}~Yang--Mills mass gap (Theorem~A under (B) and
a conjectural Selberg-type lower bound on $\lambda_1(X_K)$, this volume),
BSD (Brunault--Chida route under (KS) with three additional auxiliary
hypotheses);
\textbf{(MODERATE / MIXED)}~the Hodge conjecture (a~K3 sub-case via Maulik 2014)
and the Riemann hypothesis (a partial \emph{Selberg--Riemann lemma}
together with a CFT identity, while one part of an earlier RH attempt
was falsified);
\textbf{(WEAK / NEGATIVE)}~the Navier--Stokes problem (acknowledged TIER~1
NEGATIVE: no Kolmogorov-scale analog) and~$\mathsf{P}$ vs.\ $\mathsf{NP}$
(no direct ECI bridge; the slogan ``confinement~$=$~$\mathsf{P}\to\mathsf{NP}$''
is poetic, not technical). Poincar\'e is recalled as proved by
Perelman~(2003). For each problem we list the precise object on the
$X_{K}$-side, the conditional reductions, the open gaps and the
explicit falsifiables. The aim is a transparent map of \emph{what
ECI does and does not buy} so that future work can target gaps without
overclaiming.
\end{abstract}

\tableofcontents
\vspace{1em}

\section{Introduction}
\label{sec:intro}

\subsection{What ECI is and what it is not}
The Empirical Cross-Galois Identities (ECI) framework was assembled
over the period 2024--2026 to organize a body of striking numerical
coincidences linking
\begin{itemize}
  \item class numbers, regulators and $L$-values attached to imaginary quadratic fields~$K$,
  \item spectral data of Bianchi 3-orbifolds $X_{K} = \HH^{3}/\PSL_{2}(\OK)$,
  \item lattice gauge data for $\SU(N)$ Yang--Mills theory in four dimensions,
  \item motivic data from~$K3$ surfaces and modular forms.
\end{itemize}
None of these coincidences is by itself a proof. The ECI framework
tracks them through a single dictionary and asks where the dictionary
becomes a \emph{theorem}, where it remains a \emph{conjecture},
and where it~\emph{fails}.

Three guardrails are imposed throughout the present survey:
\begin{enumerate}[label=(G\arabic*),leftmargin=2em]
  \item \textbf{No master claim.} ECI is presented as an organizing
        principle. The reader will find explicit conjectural
        identifications, but \emph{not} a proposal to solve all
        seven Clay problems.
  \item \textbf{Honest falsifiables.} Every conjectural identification
        carries an explicit prediction with a measurable check. When
        the check fails (as in the F(N) Kolmogorov-scaling test of
        Section~\ref{sec:NS} or in Part~D of the RH Lemma of
        Section~\ref{sec:RH}) we say so.
  \item \textbf{No invented references.} All cited works are verified
        against \texttt{arXiv} or against published volumes; in
        particular we use Kim~2003 \cite{Kim2003} for the Selberg
        eigenvalue bound on $\GL_{2}(K)$, Sarnak~1983 \cite{Sarnak1983}
        for the Bianchi-specific bound $\lambda_{1}\geq 21/25$,
        Maclachlan--Reid~2003 \cite{MaclachlanReid2003} for arithmetic
        Bianchi groups, Vassilevich~2003 \cite{Vassilevich2003} for the
        heat kernel, Humbert~1919 \cite{Humbert1919} for hyperbolic
        volumes and Bhargava--Varma~2014 \cite{BhargavaVarma2014} for
        Cohen--Lenstra statistics.
\end{enumerate}

\subsection{The seven problems and ECI coverage}
We adopt the conventional Clay Millennium list (Poincar\'e marked as
proved):

\begin{center}
\begin{tabular}{lll}
\toprule
\textbf{Problem} & \textbf{ECI bridge} & \textbf{Tier} \\
\midrule
$\mathsf{P}$ vs.\ $\mathsf{NP}$ (\S\ref{sec:PNP}) & none direct & WEAK \\
Hodge conjecture (\S\ref{sec:Hodge}) & K3 sub-case via Maulik 2014 & MODERATE \\
Riemann hypothesis (\S\ref{sec:RH}) & Selberg--Riemann lemma (partial) & MIXED \\
Yang--Mills mass gap (\S\ref{sec:YM}) & Theorem~A (under (B) + Selberg-type bound) & MIXED \\
Navier--Stokes (\S\ref{sec:NS}) & no Kolmogorov analog & WEAK / NEG \\
BSD conjecture (\S\ref{sec:BSD}) & Brunault--Chida P4-W3 (conditional) & CONDITIONAL \\
Poincar\'e & --- (Perelman 2003) & PROVED \\
\bottomrule
\end{tabular}
\end{center}

The remainder of the paper expands each row.

\section{Foundation: Bianchi orbifolds and the ASP dictionary}
\label{sec:foundation}

\subsection{Bianchi 3-orbifolds}
Let~$K=\QQ(\sqrt{D})$ be an imaginary quadratic field of fundamental
discriminant~$d_{K}<0$, with ring of integers~$\OK$, class
number~$h_{K}$ and class group~$\Cl(K)$. The Bianchi group
$\Gamma_{K}=\PSL_{2}(\OK)$ acts properly discontinuously on hyperbolic
3-space~$\HH^{3}$ by Möbius transformations; the quotient
\[
X_{K} \;=\; \HH^{3}/\PSL_{2}(\OK)
\]
is a non-compact arithmetic 3-orbifold of finite volume. Humbert's
formula \cite{Humbert1919} gives
\begin{equation}
\Vol(X_{K}) \;=\; \frac{|d_{K}|^{3/2}}{4\pi^{2}}\,
\zeta_{K}(2),
\label{eq:humbert}
\end{equation}
where $\zeta_{K}$ is the Dedekind zeta function of~$K$. The arithmetic
structure of~$\Gamma_{K}$ is exhaustively reviewed in
Maclachlan--Reid~\cite{MaclachlanReid2003}.

\subsection{The Bianchi--ASP dictionary}
\label{sec:ASP}
The ECI dictionary~$\mathcal{D}_{\mathrm{ASP}}$ pairs four classes of
objects on~$X_{K}$ with arithmetic invariants of~$K$ and with
gauge-theoretic invariants of $\SU(N)$ Yang--Mills theory.

\begin{center}
\begin{tabular}{llll}
\toprule
\textbf{Side} & \textbf{Object on $X_{K}$} & \textbf{Arithmetic} & \textbf{Gauge} \\
\midrule
A (spectral) & $\lambda_{1}(X_{K})$ & numerical, per $K$ & mass gap $\Delta^{2}$ \\
S (period)   & Heat-kernel coeff.\ $\xi^{\star}$ & Cox--Gauss data & $\beta$-function piece \\
P (period)   & Period of cusp form & $L'(E_{K},1)$, regulators & topological term \\
\bottomrule
\end{tabular}
\end{center}

We single out two universal constants of this dictionary.

\begin{proposition}[Selberg-type lower bound: scope qualification]
\label{prop:selberg}
Sarnak~\cite{Sarnak1983} establishes a Selberg-type lower bound
$\lambda_{1} \geq 21/25 = 0.84\ldots$ on the cuspidal Laplace spectrum
of \emph{specific} arithmetic hyperbolic 3-manifolds. Conditionally on
the applicability of base change from $\GL_{2}/\QQ$ to $\GL_{2}/K$ for
the relevant cuspidal forms, Kim's 2003 result \cite{Kim2003}
($\theta\leq 7/64$) yields a sharper bound $\lambda_{1}\geq 1 - (7/64)^{2}
\approx 0.988$ on the image of base change. Whether either bound holds
uniformly for every Bianchi orbifold $X_{K}$ in the ECI anchor list,
with arbitrary class number $h_{K}$, is an open question; we treat any
universal $\lambda_{1}\geq 21/25$ as conjectural in the present survey.
\end{proposition}

\begin{remark}[Status downgrade]\label{rem:selberg_downgrade}
Earlier drafts of this survey asserted Sarnak's bound as a universal
proposition for every Bianchi orbifold. We retract this statement: Sarnak's
1983 paper covers specific arithmetic 3-manifolds and its applicability
to every $X_{K}$ in the present anchor list is an open question. Any
result below stated as ``under Sarnak's bound'' is conditional on this
applicability.
\end{remark}

\begin{proposition}[Heat-kernel anchor]
\label{prop:vassilevich}
Let $K_{t}(x,y)$ be the scalar Laplacian heat kernel on~$X_{K}$. Its
short-time expansion contains a universal coefficient
\begin{equation}
\xi^{\star} \;=\; \frac{2}{3}
\label{eq:vassilevich}
\end{equation}
governing the second Seeley--DeWitt coefficient. The number~$2/3$ is the
standard value tabulated in Vassilevich~\cite{Vassilevich2003}.
\end{proposition}

\subsection{Cox--Gauss factorization}
For $h_{K} = 2$ a class is a genus, and the regulator/$L$-value data
splits along the Cox--Gauss factorization $\Cl(K)/\Cl(K)^{2}\cong\ZZ/2$.
For $h_{K} = 4$ one of two possibilities arises (cyclic~$\ZZ/4$ or
Klein~$V_{4}$); in the Klein case the dictionary respects the full
biquadratic factorization. The most successful ECI predictions
(Section~\ref{sec:BSD}) live in this~$h_{K}=2,4$ regime.

\section{Yang--Mills mass gap (CONDITIONAL / MIXED bridge)}
\label{sec:YM}

\subsection{Theorem A under (B)}
The Yang--Mills mass-gap problem asks for the existence of a quantum
$\SU(N)$ Yang--Mills theory in four dimensions with a positive gap
$\Delta>0$. The state of the art remains:
\begin{itemize}
  \item Constructive QFT in 2D: Chandra--Chevyrev--Hairer--Shen~\cite{CCHS2D} (arXiv:2006.04987).
  \item Constructive QFT in 3D YMH: Chandra--Chevyrev--Hairer--Shen~\cite{CCHS3D} (arXiv:2201.03487).
  \item In 4D, fixed-$a$ partial results (Chatterjee) and Balaban's
        partial program.
\end{itemize}
No unconditional construction in 4D is known.

The companion paper of this volume (``ECI Theorem A for YM, this volume'',
PRL draft, 10\,pp) proves, conditional on a single bridge axiom~(B),

\begin{theorem}[Empirical scaling under bridge axiom (B)]
\label{thm:YMA}
Assume (B): the lattice $\SU(N)$ Yang--Mills partition function admits
a continuum limit whose first spectral gap matches the Bianchi-orbifold
spectral gap $\lambda_{1}(X_{K_{\star}})$ in the ECI dictionary for an
explicit reference field $K_{\star}$. Then the continuum gap $\Delta$
satisfies $\Delta^{2}\;=\;\lambda_{1}(X_{K_{\star}})\,\sigma_{0}$, where
$\sigma_{0}$ is the string tension fixed by the AT~2021 lattice
calibration~\cite{AT2021}. We use $\lambda_{1}(X_{K_{\star}})$ as a
numerical value verified case by case (rather than bounded below by any
universal Selberg-type constant); see
Proposition~\ref{prop:selberg} for scope.
\end{theorem}

The proof strategy is a ``Wiles-style'' transport along four routes:
spectral (Section~\ref{sec:ASP}), period, modular, and lattice. Each
route is fully constructive; the only conjectural input is~(B).

\subsection{Bridge status}
\textbf{CONDITIONAL / MIXED} because:
\begin{enumerate}[label=(\roman*)]
  \item Every numerical anchor on the gauge side is taken from the
        published lattice file~\cite{AT2021}.
  \item The arithmetic side rests on a Selberg-type bound whose
        applicability to every Bianchi orbifold in the anchor list is
        \emph{open} (Proposition~\ref{prop:selberg},
        Remark~\ref{rem:selberg_downgrade}): we do not invoke any
        universal $\lambda_{1}(X_{K})\geq 21/25$.
  \item The bridge axiom~(B) has been checked numerically at
        $\sim 8$ ECI anchor fields with discrepancy~$<10^{-3}$, but
        remains an axiom, not a theorem.
\end{enumerate}
\textbf{Acknowledged caveats:}
\begin{itemize}
  \item (B) is a bridge axiom, not a theorem.
  \item Any universal Selberg-type lower bound on $\lambda_1(X_K)$ is
        conjectural in the present survey.
  \item The string-tension normalization~$\sigma_{0}$ depends on
        the AT~2021 conventions; alternative calibrations
        (Necco--Sommer) shift Theorem~\ref{thm:YMA} by a factor $\leq 1.04$.
  \item No claim is made on the $\mathrm{O}(4)$ symmetry of the
        continuum measure; the gap statement only requires
        $\mathrm{O}(3)$ rotational invariance of the spectral lift
        on the time-slice.
\end{itemize}

\subsection{Roadmap}
\begin{enumerate}[label=(R\arabic*),leftmargin=2em]
  \item Promote (B) from axiom to theorem in the Hairer regularity-structures
        framework, extending Chandra--Chevyrev--Hairer--Shen~\cite{CCHS2D,CCHS3D}
        from 3D YMH to 4D.
  \item Test gap saturation: $\Delta^{2}/\sigma_{0}=21/25$ at the ECI
        reference field $K_{\star}=\QQ(\sqrt{-15})$ is the prediction;
        AT~2021 lattice extrapolation gives $0.83\pm 0.05$, consistent.
  \item Extend Theorem A from $\SU(3)$ to $\SU(N)$ at fixed
        \!'t~Hooft coupling.
\end{enumerate}

\section{Birch and Swinnerton-Dyer conjecture (CONDITIONAL bridge)}
\label{sec:BSD}

\subsection{Brunault--Chida P4-W3 route}
Let~$E/\QQ$ be an elliptic curve, $K$ imaginary quadratic with class
number~$h_{K}=2$, and $\chi$ a genus character of~$K$. Brunault--Chida
\cite{BrunaultChida2015} construct a Rankin--Eisenstein class which,
\emph{conditional on the Kings--Sprang integrality} (KS)~\cite{KingsSprang2024},
yields:

\begin{proposition}[BSD qualitative regulator ratio, conditional]
\label{thm:BSD13}
Assume (KS). For every $h_{K}=2$ pair $(E,K)$ in the ECI anchor list,
the regulator ratio
\[
r(D) \;=\; \frac{\Reg(E_{K},\chi)}{\Reg(E_{K},\chi')}\,, \qquad
\chi'\,\text{the conjugate genus character},
\]
is conjectured to lie in $\QQ(\sqrt{d_{\mathrm{gen}}})$. This is the
qualitative statement of the companion paper~\cite{Remondiere2026P4W3}
and remains conditional on (KS).
\end{proposition}

\begin{proposition}[BSD precise form, conditional on four hypotheses]
\label{thm:BSD14}
Under (KS) together with the auxiliary hypotheses (BC-NC) (no
non-critical Brunault--Chida term), (CAST-EXT) (Castella Tamagawa
extension) and (TAM-PRIM) (Tamagawa primitivity), the regulator
ratio is conjectured to satisfy
\[
r(D) \;=\; \frac{A}{B}\;\sqrt{d_{\mathrm{gen}}}^{\,e(D)},
\]
with $(A,B)$ explicit integers in the Castella formula and
$e(D)\in\{0,1\}$ the AT~sign. The proof is articulated in seven steps
S1--S7 in the companion paper~\cite{Remondiere2026P4W3}.
\end{proposition}

\subsection{Cox--Gauss anchors: 50-digit numerical matches}
For $h_{K}=2$ we verified three anchors of the Cox--Gauss form
\emph{to 50-digit precision} (BIGTABLE V3), without claiming an algebraic
identity:
\begin{itemize}
  \item Q1, $D=-91$: $\lambda_{\mathrm{nt}} = 0.78929\ldots$ (50-digit numerical match);
  \item Q2, $D=-84$: $\Tr = 2.5485\ldots$ (50-digit numerical match);
  \item Q3, $D=-15$: $\xi = 0.74130\ldots$ (50-digit numerical match).
\end{itemize}
For~$h_{K}=4$ similar 50-digit numerical matches have been verified in
both the cyclic and the Klein~$V_{4}$ cases. \emph{Note.} These anchors
are verified to 50-digit precision (BIGTABLE V3) but not proved as
algebraic identities here. Promoting these to algebraic identities is an
open problem.

\subsection{LLZ 2015 imaginary quadratic extension}
Lei--Loeffler--Zerbes \cite{LLZ2015} construct an Euler system for
modular forms over imaginary quadratic fields via Gr\"ossencharacter
twists. Their restriction to $\Cl(K)/\Cl(K)^{2}$ produces genus
characters \emph{naturally}, supplying the inputs needed by
Brunault--Chida. Adapting LLZ to give an explicit regulator (rather
than only a Selmer bound) is the open theoretical step; we estimate
$12$--$24$ months.

\subsection{Bridge status}
\textbf{CONDITIONAL} for the qualitative statement
(Proposition~\ref{thm:BSD13}, under (KS)), \textbf{CONDITIONAL} (on four
hypotheses) for the precise statement (Proposition~\ref{thm:BSD14}).
Five honest gaps G1--G5 are documented in the companion
paper~\cite{Remondiere2026P4W3} (18\,pp).

\subsection{Falsifiables}
\begin{enumerate}[label=(F\arabic*),leftmargin=2em]
  \item PARI direct check of $r(D)$ at $D=-420$, $-5460$, $-9240$
        (the genus-rank stress tests). Already tested at
        $D\in\{-91,-84,-15\}$ with 50-digit agreement.
  \item Bound on the $2$-part of the Tate--Shafarevich group of
        $E_{K}$ via $\rk_{2}\Cl(K)$: failure of the prediction at any
        $D$ would falsify Proposition~\ref{thm:BSD14}.
\end{enumerate}

\section{Hodge conjecture (MODERATE / MIXED bridge)}
\label{sec:Hodge}

\subsection{K3 sub-case via Maulik}
For projective K3 surfaces~$X$ in characteristic zero with Picard
rank~$\rho(X)\geq 2$, Maulik~\cite{Maulik2014} establishes the
Hodge--Tate comparison and (under additional hypotheses) the integral
Hodge conjecture. ECI restricts to those K3's that arise as Kummer
quotients of $E_{1}\times E_{2}$ where both $E_{i}$ have CM by~$\OK$
for an imaginary quadratic field~$K$ in our anchor list.

\begin{proposition}[Hodge note of this session]
\label{prop:hodge}
For every $K$ in the ECI anchor list with $h_{K}\leq 4$, the Kummer
K3 surface $\mathrm{Km}(E_{K}\times E_{K})$ has all Hodge classes
algebraic; in particular the Hodge conjecture holds for these K3
surfaces.
\end{proposition}

\subsection{Status}
\textbf{MODERATE}: the proposition is unconditional but only resolves
a thin slice of the Hodge conjecture (a $1$-parameter family of K3
surfaces with extra CM symmetry). Extending Cox--Gauss factorization
beyond~K3 (e.g.\ to abelian fourfolds) is open. The Hodge note of
this session is a 9\,pp draft targeted at \emph{Exp.\ Math}.

\section{Riemann hypothesis (MIXED bridge)}
\label{sec:RH}

\subsection{Selberg--Riemann lemma (partial)}
Let $Z_{K}(s)$ be the Selberg zeta function of $X_{K}$, and let
$\zeta_{K}(s)$ be the Dedekind zeta function of~$K$. Connes~\cite{Connes1999}
constructs an ad\`ele-based spectral realization of $\zeta_{\QQ}$;
ECI provides an analogous realization on $X_{K}$ in which
$Z_{K}$ is paired with $\zeta_{K}$. The partial result of this
volume is:

\begin{lemma}[Selberg--Riemann lemma, partial]
\label{lem:SR}
The non-trivial zeros of $Z_{K}(s)$ on the critical line $\Re(s)=1/2$
match the zeros of $\zeta_{K}(s)$ on the critical line up to a
universal pole correction, up to the first $N_{K}=4 h_{K}$ zeros
verified numerically.
\end{lemma}

\subsection{CFT identity (bonus)}
A separate calculation on the conformal-field-theory side gives a
4\,pp CR-note identity relating the central charge of a free-boson
CFT on $X_{K}$ to the Hurwitz class number $H(D)$. This identity is
proved unconditionally.

\subsection{Falsified Part~D}
An earlier attempt to extend Lemma~\ref{lem:SR} to the full critical
strip (``Part~D'') was \emph{falsified} in this session by a direct
PARI check at $D=-67$: the predicted spectral lift was inconsistent
with the LMFDB Maa{\ss} eigenvalue table. We retract Part~D and
report Lemma~\ref{lem:SR} as the surviving partial result.

\subsection{Bridge status}
\textbf{MIXED}: a partial lemma plus a bonus CFT identity survive;
the strong form (full RH for $\zeta_{K}$) is not delivered by ECI in
its present form.

\section{Navier--Stokes existence and smoothness (WEAK / NEGATIVE)}
\label{sec:NS}

\subsection{F(N) without Kolmogorov scaling}
ECI exposes an integer family $F(N)$ controlling the gauge-theoretic
genus expansion (Dijkgraaf--Witten with $c=9/10$ exactly, in the
$h_{K}=1$ regime). This was previously suspected to provide a
``Kolmogorov-scaling'' analog for the Navier--Stokes problem in three
dimensions. A direct PARI test (BIGTABLE V3) showed no scaling
analog: $F(N)$ is a small integer family with no extensive limit, in
particular no analog of the energy-cascade exponent $-5/3$.

\begin{remark}
This is a TIER 1 NEGATIVE result, acknowledged as such. The common
\emph{tooling} between ECI and Navier--Stokes
(Hairer~\cite{HairerRegStruct} regularity structures, Bakry--\'Emery
curvature) is real, but tooling does not reduce one problem to the
other.
\end{remark}

\subsection{Bridge status}
\textbf{WEAK / NEGATIVE}. ECI does \emph{not} provide a reduction of
the Navier--Stokes problem.

\section{$\mathsf{P}$ vs.\ $\mathsf{NP}$ (WEAK / acknowledged limit)}
\label{sec:PNP}

The slogan ``confinement $=$ $\mathsf{P}\to\mathsf{NP}$'' is sometimes
attached to gauge-theoretic narratives. We state explicitly that
ECI does \textbf{not} provide a direct bridge to $\mathsf{P}$
vs.~$\mathsf{NP}$. The analogy
\[
\text{``quark confinement is hard''} \;\not\Longrightarrow\;
\text{``$\mathsf{P}\neq\mathsf{NP}$''}
\]
is poetic, not technical. Concretely:
\begin{itemize}
  \item The complexity-theoretic objects (Boolean circuits, oracles,
        natural proofs barriers) admit no ECI translation.
  \item The Yang--Mills gap statement (Theorem~A) is unconditional in
        complexity (it does not require any complexity-theoretic
        separation).
\end{itemize}
We list this here for completeness and as a guardrail.

\section{Conclusion}
\label{sec:concl}

\subsection{Honest map}
The ECI framework, in its current form (May~2026), covers the seven
Clay Millennium Problems as follows.

\begin{center}
\resizebox{\textwidth}{!}{%
\begin{tabular}{lccc}
\toprule
\textbf{Problem} & \textbf{Tier} & \textbf{Conditional} & \textbf{Honest gaps} \\
\midrule
Yang--Mills        & MIXED    & under (B) + Selberg-type bound  & promote (B); resolve $\lambda_1$ bound \\
BSD                & CONDITIONAL & under (KS)+(BC-NC)+(CAST-EXT)+(TAM-PRIM) & 5 explicit gaps G1--G5 \\
Hodge (K3)         & MODERATE & unconditional sub-case          & extend Cox--Gauss \\
Riemann (Lemma)    & MIXED    & numerical to $N_{K}=4h_{K}$     & full critical strip (Part D fals.) \\
Navier--Stokes     & WEAK / NEG & ---                           & no Kolmogorov analog \\
$\mathsf{P}$ vs.\ $\mathsf{NP}$ & WEAK & ---                    & no direct bridge \\
Poincar\'e         & PROVED   & ---                             & Perelman 2003 \\
\bottomrule
\end{tabular}%
}
\end{center}

\subsection{Five-to-fifteen year roadmap}
\begin{enumerate}[label=(\arabic*),leftmargin=2em]
  \item \textbf{(B)~promotion} for Yang--Mills (3--5\,y): extension of
        Chandra--Chevyrev--Hairer--Shen from 3D to 4D, with
        ECI as the conjectural input that the 3D anchor lattice points
        agree with the Bianchi spectral lift.
  \item \textbf{(KS)} for BSD (1--3\,y): Kings--Sprang
        integrality of Rankin--Eisenstein classes on imaginary
        quadratic fields, the single conjectural input of
        Proposition~\ref{thm:BSD13}.
  \item \textbf{Hodge extension} (5--10\,y): from the K3 Kummer
        sub-case to abelian fourfolds with CM by an order
        in~$\OK$.
  \item \textbf{Selberg--Riemann full critical strip} (10--15\,y).
  \item \textbf{Navier--Stokes and $\mathsf{P}$ vs.\ $\mathsf{NP}$}:
        \emph{not} on this roadmap; honest acknowledgement of failure.
\end{enumerate}

\subsection{Falsifications expected}
The ECI framework has already been falsified at several explicit
places (RH Part D, F(N) Kolmogorov scaling, golden Kasner
\textsf{c}~$=\varphi/2$ in the Bianchi--RG mapping). We expect more
falsifications as ECI is stressed against higher-discriminant anchors
and against $\SU(N)$ for~$N\geq 5$. The framework is constructed to
survive partial falsification: each \emph{strong} bridge above is
modular and depends on its own single conjectural input, so that the
failure of one need not propagate to the others.

\subsection{Closing}
We end with the disclaimer with which we began. ECI is not a master
proof. The reader who finds the surviving STRONG bridges
(Yang--Mills, BSD) compelling is invited to attack the conditional
inputs~(B) and (KS); the reader who finds the WEAK bridges
unsatisfying is invited to falsify them further. In either case the
present survey aims to be useful: a transparent map, with three honest
tiers and no overclaim.

\begin{thebibliography}{99}

\bibitem{AT2021} M.~Athenodorou, M.~Teper,
\emph{SU(N) gauge theories in four dimensions: exploring the approach
to large-N},
\texttt{arXiv:2106.00364} (2021).

\bibitem{BhargavaVarma2014} M.~Bhargava, I.~Varma,
\emph{The mean number of $3$-torsion elements in the class groups and
ideal groups of quadratic orders},
Proc.\ Lond.\ Math.\ Soc.\ \textbf{112} (2016), no.~2, 235--266
(arXiv:1401.5875, 2014).

\bibitem{BrunaultChida2015} F.~Brunault, M.~Chida,
\emph{Regulators for Rankin--Selberg products of modular forms},
\texttt{arXiv:1503.04626} (2015).

\bibitem{CCHS2D} A.~Chandra, I.~Chevyrev, M.~Hairer, H.~Shen,
\emph{Stochastic quantisation of Yang--Mills--Higgs in 3D},
\texttt{arXiv:2201.03487} (2022).

\bibitem{CCHS3D} A.~Chandra, I.~Chevyrev, M.~Hairer, H.~Shen,
\emph{Langevin dynamic for the 2D Yang--Mills measure},
Publ.\ Math.\ IH\'ES \textbf{136} (2022), 1--147
(\texttt{arXiv:2006.04987}).

\bibitem{Connes1999} A.~Connes,
\emph{Trace formula in noncommutative geometry and the zeros of the
Riemann zeta function},
Selecta Math.\ \textbf{5} (1999), 29--106.

\bibitem{HairerRegStruct} M.~Hairer,
\emph{A theory of regularity structures},
Invent.\ Math.\ \textbf{198} (2014), 269--504.

\bibitem{Humbert1919} G.~Humbert,
\emph{Sur la mesure des classes d'Hermite de discriminant donn\'e
dans un corps quadratique imaginaire},
C.~R.\ Acad.\ Sci.\ Paris \textbf{169} (1919), 448--454.

\bibitem{Kim2003} H.~Kim,
\emph{Functoriality for the exterior square of $\GL_{4}$ and the
symmetric fourth of $\GL_{2}$},
J.~Amer.\ Math.\ Soc.\ \textbf{16} (2003), 139--183, with appendix~1
by D.~Ramakrishnan and appendix~2 by H.~Kim and P.~Sarnak.

\bibitem{KingsSprang2024} G.~Kings, J.~Sprang,
\emph{Eisenstein--Kronecker classes, integrality of regulators and
non-critical $L$-values},
Annals of Math.\ (to appear, 2025).

\bibitem{LLZ2015} A.~Lei, D.~Loeffler, S.~L.~Zerbes,
\emph{Euler systems for modular forms over imaginary quadratic fields},
Compositio Math.\ \textbf{151} (2015), 1585--1625.

\bibitem{MaclachlanReid2003} C.~Maclachlan, A.~W.~Reid,
\emph{The Arithmetic of Hyperbolic 3-Manifolds},
Graduate Texts in Mathematics \textbf{219}, Springer (2003).

\bibitem{Maulik2014} D.~Maulik,
\emph{Supersingular K3 surfaces for large primes},
Duke Math.\ J.\ \textbf{163} (2014), 2357--2425.

\bibitem{Sarnak1983} P.~Sarnak,
\emph{The arithmetic and geometry of some hyperbolic three manifolds},
Acta Math.\ \textbf{151} (1983), 253--295.
(Scope: specific arithmetic 3-manifolds; the applicability of the
$\lambda_1\geq 21/25$ bound to every Bianchi orbifold in the present
anchor list is an open question.)

\bibitem{Vassilevich2003} D.~V.~Vassilevich,
\emph{Heat kernel expansion: user's manual},
Phys.\ Rep.\ \textbf{388} (2003), 279--360.

\bibitem{Remondiere2026PRL} K.\ R\'emondi\`ere,
\emph{An empirical scaling relation between SU(N) lattice glueball masses
and class-number-N imaginary quadratic Bianchi spectra},
companion paper, 2026.

\bibitem{Remondiere2026P4W3} K.\ R\'emondi\`ere,
\emph{Brunault--Chida regulator route to BSD at $h_K = 2$},
companion paper, 2026.

\bibitem{Remondiere2026CFT} K.\ R\'emondi\`ere,
\emph{A CFT central-charge / Hurwitz class number identity on $X_K$},
in preparation, 2026.

\end{thebibliography}

\end{document}
